\subsection{مبانی نظری و فرمول‌بندی ریاضی}
فرآیند فشرده‌سازی این الگوریتم در سه گام ریاضی پیوسته انجام می‌گیرد:
\begin{enumerate}
    \item \textbf{چرخش یکامتعامد مختصات (\lr{Orthogonal Coordinate Rotation}):}
    جهت همگن‌سازی توزیع مقادیر پرت فعال‌سازی‌ها در تمامی جهات، یک ماتریس یکامتعامد تصادفی هاار $\mathbf{Q} \in \mathbb{R}^{d_{\text{in}} \times d_{\text{in}}}$ با شرط $\mathbf{Q}^\top \mathbf{Q} = \mathbf{I}$ با استفاده از تجزیه \lr{QR} تولید می‌گردد. با توجه به نامتغیر بودن ضرب ماتریسی تحت نگاشت یکامتعامد:
    \begin{equation}
    \mathbf{W} \mathbf{X} = (\mathbf{W} \mathbf{Q}) (\mathbf{Q}^\top \mathbf{X}) = \tilde{\mathbf{W}} \tilde{\mathbf{X}}
    \end{equation}
    که در آن $\tilde{\mathbf{W}} = \mathbf{W} \mathbf{Q}$ و $\tilde{\mathbf{X}} = \mathbf{Q}^\top \mathbf{X}$ مقادیر وزن و فعال‌سازی در فضای مختصات چرخیده‌شده هستند.

    \item \textbf{برش زیرفضا بر پایه واریانس فعال‌سازی‌ها (\lr{Variance-Based Subspace Slicing}):}
    واریانس کانال‌های فعال‌سازی چرخیده‌شده محاسبه گردیده و $r = \lfloor d_{\text{in}} \cdot \rho \rfloor$ کانال با بیشترین واریانس به عنوان زیرفضای بااهمیت انتخاب می‌شوند:
    \begin{equation}
    \mathcal{I} = \operatorname{Top-}r(\operatorname{Var}(\tilde{\mathbf{X}})), \quad \mathbf{P}_r = [\mathbf{e}_i]_{i \in \mathcal{I}} \in \mathbb{R}^{d_{\text{in}} \times r}
    \end{equation}
    ماتریس‌های وزن و ورودی برش‌یافته به صورت زیر حاصل می‌گردند:
    \begin{equation}
    \mathbf{W}_{\text{sliced}} = \tilde{\mathbf{W}} \mathbf{P}_r \in \mathbb{R}^{d_{\text{out}} \times r}, \quad \mathbf{X}_{\text{sliced}} = \mathbf{P}_r^\top \tilde{\mathbf{X}} \in \mathbb{R}^{r \times N}
    \end{equation}

    \item \textbf{کوانتیزاسیون غیریکنواخت ۴ بیتی با جبران هسیان مرتبه دوم (\lr{Second-Order OBS INT4}):}
    ماتریس هسیان تجربی به همراه جمله میراسازی تیخونوف $\lambda$ تشکیل می‌شود:
    \begin{equation}
    \mathbf{H} = \frac{2}{N} \mathbf{X}_{\text{sliced}} \mathbf{X}_{\text{sliced}}^\top + \lambda \mathbf{I}_r
    \end{equation}
    سپس عناصر ستون‌های وزن به ترتیب با کتاب‌کد گوسی \lr{AstroQuantizer} کوانتیزه شده و خطای گردکردن به صورت معکوس روی ستون‌های بعدی جبران می‌گردد:
    \begin{equation}
    \mathbf{W}_{:, j+1:r} \leftarrow \mathbf{W}_{:, j+1:r} - \left( \mathbf{w}_j - q(\mathbf{w}_j) \right) \frac{[\mathbf{H}^{-1}]_{j, j+1:r}}{[\mathbf{H}^{-1}]_{jj}}
    \end{equation}
    که در آن $q(\cdot)$ نگاشت به نزدیک‌ترین مرکز خوشه غیرخطی در کتاب‌کد بهینه‌شده توزیع نرمال است.
\end{enumerate}

که در آن $\sigma_i(\mathbf{W})$ بیانگر $i$-امین مقدار منفرد مرتب‌شده ماتریس به صورت نزولی است. نسبت حفظ انرژی طیفی $\eta_r = \frac{\sum_{i=1}^r \sigma_i^2}{\sum_{i=1}^{d} \sigma_i^2}$ تضمین می‌کند که با انتخاب آستانه $\tau \ge 0.90$، بیش از نود درصد از پراکندگی واریانس داده‌ها حفظ شده و افت سیگنال در انتشار گرادیان رو به عقب به حداقل ممکن کاهش می‌یابد.

\paragraph{تحلیل مجانبی و پیچیدگی اشتراک حافظه سخت‌افزاری}
از منظر پیچیدگی محاسباتی و انتقال داده‌ها در سلسومراتب حافظه (حافظه سراسری \lr{HBM} به حافظه محلی \lr{SRAM})، عملیات ضرب ماتریسی در این مکانیزم نیازمند تخصیص بافر موقت با ابعاد فاکتورگیری‌شده است. برای پردازش دنباله‌ای به طول $N$ با بعد مخفی $D$، بار محاسباتی هسته از فرم درجه دو به رابطه زیر تقلیل می‌یابد:
\begin{equation}
\label{eq:complexity_rotational_slicegpt_int4}
\mathcal{C}_{\text{rotational\_slicegpt\_int4}} = \mathcal{O}\left( B \cdot N \cdot D \cdot r + B \cdot H \cdot N \cdot d_k \right) \ll \mathcal{O}\left( B \cdot H \cdot N^2 \cdot d_k \right)
\end{equation}
این بهبود مرتبه پیچیدگی، با کاهش چشمگیر حجم تبادل تانسورها در گذرگاه‌های ارتباطی کارت گرافیک، مانع از اشباع پهنای باند حافظه (پدیده \lr{Memory-Bound}) در فرآیند استنتاج طولانی‌مدت می‌گردد.
